% ===================================================================== Chapter 15
\chapter{Parametric Study}\label{ch:param}

\section{Study Design}
The parametric study asks how the flow, thermal and structural responses change with the two operating variables and one design variable
(Table~\ref{tab:pmatrix}). It is a one-factor-at-a-time design around P00 with the S1 supports throughout; the support question is treated
separately (Chapter~\ref{ch:supports}). Every case is an actual simulation chain:
\begin{itemize}
\item a converged Fluent CFD solution of its own (Section 9B-1), with the same physics, staging and convergence criteria as P00;
\item mapping of its own solid temperature field with the method of Chapter~\ref{ch:mapping};
\item Mechanical LC1 and LC2 solutions and a linear buckling solution on the LC2 pre-stress (Section 9B-2).
\end{itemize}
No screening estimate, uniform temperature or interpolated value enters any result of this chapter. A control case, C00, repeated the whole
pipeline for P00 and reproduced it (LC2 peak within $9.5\times10^{-6}$, $\lambda_1$ within $4.2\times10^{-7}$).

%% generated by gen_tables.py - RE-ANALYSIS 2026
\begin{table}[htbp]
\centering
\footnotesize
\caption{Parametric case matrix (one factor at a time around P00; S1 supports)}
\label{tab:pmatrix}
\begin{tabular}{lll>{\raggedright\arraybackslash}p{0.38\linewidth}}
\toprule
Case & Variable & Value (change) & Held constant / mesh \\
\midrule
P00 & -- & baseline & 159,840 CFD cells; mesh B \\
V01 / V03 & inlet velocity & 21.15 / 25.85 m/s ($\mp$10\,\%) & $q''$, geometry; P00 meshes \\
Q01 / Q03 & outer-wall heat flux & 7200 / 8800 W/m$^2$ ($\mp$10\,\%) & $V$, geometry; P00 meshes \\
T01 / T03 & wall thickness & 8 / 12 mm ($\mp$20\,\%; $D_o$ 36 / 44 mm) & total heat input $Q$ held ($q''$ = 8888.9 / 7272.7 W/m$^2$); new CFD meshes 155,520 / 168,480 cells, FE meshes 89,424 / 127,080 nodes \\
\bottomrule
\end{tabular}
\par\smallskip\parbox{\linewidth}{\footnotesize\raggedright Source: \texttt{PROJECT\_STATE.md} \S17--\S19. Range limits: the 600~K air property table (low $V$, high $q''$) and the 128,000-node licence (thick wall).}
\end{table}


The ranges are limited by the solver data rather than chosen arbitrarily. Lower velocity and higher heat flux raise the air temperature at
the wall; the air property table ends at 600~K, and the $-$10\,\% velocity and +10\,\% heat-flux cases reach 580.4~K and 583.4~K on the air side.
The 12~mm wall is the thickest that a 2~mm radial element size allows under the 128,000-node licence (127,080~nodes).

\section{Velocity Cases}
V01 and V03 change the inlet velocity by $\mp$10\,\% at unchanged heat input. Velocity sets the cooling capacity: at 21.15~m/s the air leaves
7.6~K hotter, the maximum solid temperature rises by 24.9~K and the solid mean by 21.5~K; at 25.85~m/s the changes are $-$6.2~K, $-$20.5~K and
$-$17.7~K. The pressure drop changes by $-$13.3 and +14.1\,\% (Figure~\ref{fig:pv}a), a local exponent of about 1.37, lower than the
1.75--2 of isothermal turbulent flow because part of the static pressure drop accelerates the heated air, and that part grows when the air
heats more.

\begin{figure}[htbp]
\centering
\begin{subfigure}[t]{0.325\linewidth}\centering\includegraphics[width=\linewidth]{F01_V_vs_pressure_drop.png}\caption{Pressure drop}\end{subfigure}\hfill
\begin{subfigure}[t]{0.325\linewidth}\centering\includegraphics[width=\linewidth]{F03_V_vs_max_solid_temperature.png}\caption{Maximum solid temperature}\end{subfigure}\hfill
\begin{subfigure}[t]{0.325\linewidth}\centering\includegraphics[width=\linewidth]{F05_V_vs_lambda1.png}\caption{First buckling load factor (S1)}\end{subfigure}
\caption[Effect of inlet velocity (21.15, 23.5, 25.85~m/s)]{Effect of inlet velocity (21.15, 23.5, 25.85~m/s); the dashed line only connects the three solved cases.}
\label{fig:pv}
\src{\provdata}
\end{figure}

\section{Heat-Flux Cases}
Q01 and Q03 change the outer-wall heat flux by $\mp$10\,\%. The heat flux is the thermal load itself: the maximum solid temperature changes by
$-$28.1 and +28.6~K and the outlet temperature by $\mp$6.9~K (Figure~\ref{fig:pq}). The solid-mean temperature rise above the inlet changes by
$-$10.8 and +11.0\,\%, slightly more than in proportion; the cause (temperature-dependent air and Inconel properties are a plausible
explanation) was not isolated. The pressure drop changes by only $\pm$4.2\,\%, through the air temperature.

\begin{figure}[htbp]
\centering
\begin{subfigure}[t]{0.325\linewidth}\centering\includegraphics[width=\linewidth]{F07_q_vs_max_solid_temperature.png}\caption{Maximum solid temperature}\end{subfigure}\hfill
\begin{subfigure}[t]{0.325\linewidth}\centering\includegraphics[width=\linewidth]{F08_q_vs_LC2_stress.png}\caption{LC2 maximum von Mises stress}\end{subfigure}\hfill
\begin{subfigure}[t]{0.325\linewidth}\centering\includegraphics[width=\linewidth]{F09_q_vs_lambda1.png}\caption{First buckling load factor (S1)}\end{subfigure}
\caption{Effect of outer-wall heat flux (7200, 8000, 8800~W/m$^2$).}
\label{fig:pq}
\src{\provdata}
\end{figure}

\section{Wall-Thickness Cases}
T01 and T03 change the wall thickness to 8 and 12~mm (outer diameter 36 and 44~mm) with the bore unchanged. To isolate the section effect,
the total heat input is held at the baseline value, so the outer-wall flux becomes 8888.9 and 7272.7~W/m$^2$. New CAD models, CFD meshes
(9 and 12 solid layers) and structural meshes (four and six 2.0~mm elements through the wall) were built. The outlet temperature and pressure drop are unchanged,
and the solid temperatures move by less than 1~K: the through-wall temperature difference (6.44 / 7.58 / 8.61~K for 8 / 10 / 12~mm) and the
axial conduction through the changed cross-section act in opposite directions. A check solution of T03 with 18 instead of 12 solid layers
changed the node temperatures by at most 0.197~K.


\section{CFD Results}
Table~\ref{tab:pcfd} lists the CFD results. Every case converged at iteration 600 and held at 700, with mass and energy imbalances of at most
$4.2\times10^{-14}$ and $1.3\times10^{-12}$, $y^+_{max}$ between 0.524 and 0.645, and all temperatures inside the property tables. An independent
verification script re-derived the CFD results from the raw case outputs (309 of 309 checks passed).

%% generated by gen_tables.py - RE-ANALYSIS 2026
\begin{table}[htbp]
\centering
\footnotesize
\caption{Parametric CFD results (actual Fluent solutions, 9B-1)}
\label{tab:pcfd}
\begin{tabular}{lrrrrrrr}
\toprule
Case & $\Delta p$ [Pa] & $T_{out}$ [K] & $T_{max}$ [K] & $Q$ [W] & $Re_{in}$ & $T_{int,max}$ [K] & $\bar T_s$ [K] \\
\midrule
P00 & 438.13 & 368.93 & 562.58 & 602.76 & 29,958 & 555.27 & 525.48 \\
V01 & 379.93 & 376.56 & 587.47 & 602.76 & 26,962 & 580.38 & 546.96 \\
V03 & 499.69 & 362.69 & 542.10 & 602.76 & 32,954 & 534.61 & 507.78 \\
Q01 & 419.74 & 362.06 & 534.45 & 542.48 & 29,958 & 527.65 & 501.07 \\
Q03 & 456.64 & 375.79 & 591.13 & 663.03 & 29,958 & 583.35 & 550.31 \\
T01 & 438.11 & 368.93 & 561.97 & 602.76 & 29,958 & 555.76 & 524.54 \\
T03 & 438.14 & 368.93 & 563.07 & 602.76 & 29,958 & 554.78 & 526.38 \\
\bottomrule
\end{tabular}
\par\smallskip\parbox{\linewidth}{\footnotesize\raggedright Source: \texttt{MASTER\_PROJECT\_DATA.csv} ($\Delta p$, $T_{out}$, $T_{max}$: VERIFIED); $Q$, $Re_{in}$, interface maximum $T_{int,max}$ and solid mean $\bar T_s$ from \texttt{10\_Parametric\_Study/CFD\_Results/PARAMETRIC\_CFD\_TABLE.md} and \texttt{PROJECT\_STATE.md} \S18.3.}
\end{table}


\section{Structural Results}
The static structural results are given in Table~\ref{tab:pstruct}. For velocity and heat flux, every structural response follows the change
of the wall temperature: the LC1 growth, the LC2 end force, the LC2 stresses and the utilisation all change with normalised sensitivities
between 0.87 and 1.20 in magnitude. The LC2 peak von Mises stress changes slightly less than the mean axial stress (for V01 +9.56\,\% against
+9.72\,\%), because the temperature at the critical inlet-edge node changes less than the duct mean. For thickness the stress level barely
changes (LC2 peak $-$0.47 / +0.43\,\%), while the end force scales with the section area ($-$25.6 / +28.5\,\%). The critical locations do not move
in any case: LC2 at the outer edge of the inlet face, LC1 at the bore 6.5~mm from the inlet. No case is yield-limited; the highest
utilisation is 0.6446 (Q03).

%% generated by gen_tables.py - RE-ANALYSIS 2026
\begin{table}[htbp]
\centering
\footnotesize
\caption{Parametric static structural results (actual Mechanical solutions, 9B-2; S1 supports)}
\label{tab:pstruct}
\begin{tabular}{lrrrrrrrr}
\toprule
Case & LC1 def.\ [mm] & $\Delta L$ [mm] & LC1 VM [MPa] & LC2 VM [MPa] & $\bar\sigma_z$ [MPa] & $T_{crit}$ [K] & $S_y(T)$ [MPa] & Util. \\
\midrule
P00 & 1.8443 & 1.8409 & 24.282 & 605.161 & $-$582.440 & 437.99 & 1047.0 & 0.5780 \\
V01 & 2.0375 & 2.0339 & 24.925 & 663.001 & $-$639.053 & 450.38 & 1044.6 & 0.6347 \\
V03 & 1.6874 & 1.6842 & 23.532 & 557.544 & $-$535.853 & 427.76 & 1049.1 & 0.5315 \\
Q01 & 1.6284 & 1.6253 & 21.704 & 538.526 & $-$518.230 & 423.38 & 1050.0 & 0.5129 \\
Q03 & 2.0680 & 2.0641 & 26.649 & 673.024 & $-$647.830 & 452.79 & 1044.1 & 0.6446 \\
T01 & 1.8355 & 1.8328 & 20.872 & 602.309 & $-$580.048 & 429.52 & 1048.7 & 0.5743 \\
T03 & 1.8528 & 1.8487 & 27.206 & 607.735 & $-$584.750 & 445.36 & 1045.6 & 0.5813 \\
\bottomrule
\end{tabular}
\par\smallskip\parbox{\linewidth}{\footnotesize\raggedright Source: \texttt{MASTER\_PROJECT\_DATA.csv} (LC1 deformation, LC1/LC2 von Mises, utilisation: VERIFIED); $\Delta L$, mean axial stress $\bar\sigma_z$, critical temperature and local $S_y$ from \texttt{10\_Parametric\_Study/Results/PARAMETRIC\_STRUCTURAL\_RESULTS.csv}.}
\end{table}


\section{Buckling Results}
Table~\ref{tab:pbk} gives the buckling results. The mode is the global guided-sway shape in every case. Velocity and heat flux change
$\lambda_1$ through the end force, with the critical load itself almost constant ($P_{cr}$ changes by less than 0.8\,\%) because the section and
the modulus change little. Thickness changes $\lambda_1$ through the section: $P_{cr}$ scales with the bending stiffness ($-$36.6 / +49.3\,\%),
the end force with the area, and $\lambda_1 \propto I/(A\,\varepsilon_{th}) \propto r_g^2$, so the thin wall loses and the thick wall gains
stability margin in the idealised sense (Figure~\ref{fig:pt}).

Under S1, $\lambda_1$ falls below 1 for \textbf{Q03} (0.98834) and \textbf{T01} (0.94497). For these two cases the LC2 static stresses are
idealised pre-buckling equilibrium results: the linear stability criterion indicates loss of stability before the applied load level, and
they are not presented as physically stable operating states. \textbf{V01} ($\lambda_1$ = 1.00307) lies 0.31\,\% above 1, extremely close to the
limit; its side of 1 is inside the thermal and material uncertainty bands. In all seven S1 cases $\lambda_1$ is below the first-yield factor, so
every case is stability-first in the idealised model.

\begin{figure}[htbp]
\centering
\begin{subfigure}[t]{0.325\linewidth}\centering\includegraphics[width=\linewidth]{F11_t_vs_LC2_stress.png}\caption{LC2 maximum von Mises stress}\end{subfigure}\hfill
\begin{subfigure}[t]{0.325\linewidth}\centering\includegraphics[width=\linewidth]{F12_t_vs_lambda1.png}\caption{First buckling load factor (S1)}\end{subfigure}\hfill
\begin{subfigure}[t]{0.325\linewidth}\centering\includegraphics[width=\linewidth]{F13_t_vs_critical_buckling_load.png}\caption{Critical buckling load}\end{subfigure}
\caption{Effect of wall thickness (8, 10, 12~mm) at constant total heat input.}
\label{fig:pt}
\src{\provdata}
\end{figure}

%% generated by gen_tables.py - RE-ANALYSIS 2026
\begin{table}[htbp]
\centering
\footnotesize
\caption{Parametric linear buckling results under S1 (actual eigenvalue solutions, 9B-2)}
\label{tab:pbk}
\begin{tabular}{lrrrrll}
\toprule
Case & $N$ [kN] & $\lambda_1$ & $P_{cr}$ [kN] & First-yield factor & First reached & Note \\
\midrule
P00 & 548.94 & 1.10805 & 608.25 & 1.730 & buckling &  \\
V01 & 602.29 & 1.00307 & 604.14 & 1.575 & buckling & 0.31\,\% above 1 \\
V03 & 505.03 & 1.21087 & 611.53 & 1.882 & buckling &  \\
Q01 & 488.42 & 1.25464 & 612.79 & 1.950 & buckling &  \\
Q03 & 610.57 & 0.98834 & 603.45 & 1.551 & buckling & \textbf{$<$ 1: pre-buckling equilibrium} \\
T01 & 408.19 & 0.94497 & 385.73 & 1.741 & buckling & \textbf{$<$ 1: pre-buckling equilibrium} \\
T03 & 705.43 & 1.28702 & 907.89 & 1.720 & buckling &  \\
\bottomrule
\end{tabular}
\par\smallskip\parbox{\linewidth}{\footnotesize\raggedright Source: \texttt{MASTER\_PROJECT\_DATA.csv} ($\lambda_1$, $P_{cr}$: VERIFIED); end force $N$ and first-yield factor from \texttt{PROJECT\_STATE.md} \S19.3. Mode 1 is the global guided-sway shape in every case. For Q03 and T01 the LC2 static stresses are idealised pre-buckling equilibrium results: the linear stability criterion indicates loss of stability before the applied load level.}
\end{table}


\section{Sensitivity Trends}
Table~\ref{tab:psens} compares the three variables through normalised sensitivities, which remove the different step sizes. Velocity acts
most strongly on the pressure drop ($S$ = +1.37); heat flux most strongly on $\lambda_1$ ($S = -1.20$) and on the thermal growth; thickness most
strongly on the critical load ($S$ = +2.15) and the end force, with a negligible effect on temperatures ($|S| \le 0.009$). The normalised effect
of thickness on $\lambda_1$ (+0.77) is of the same order as those of velocity (+0.94) and heat flux ($-$1.20), although it acts through a
different mechanism. The 9A analytical screening, prepared before the simulations, agreed with the FE values of $\lambda_1$ to within
0.78\,\%, but placed V01 at or below 1; the FE result (1.0031) supersedes it, and no conclusion of this report uses a screening value.

%% generated by gen_tables.py - RE-ANALYSIS 2026
\begin{table}[htbp]
\centering
\footnotesize
\caption{Normalised sensitivities $S = (\Delta y/y_0)/(\Delta x/x_0)$ from the one-factor steps}
\label{tab:psens}
\begin{tabular}{lrrr}
\toprule
Response $y$ & Velocity $V$ & Heat flux $q''$ & Thickness $t$ \\
\midrule
Pressure drop $\Delta p$ & +1.367 & +0.421 & $\approx 0$ \\
Maximum solid temperature & $-$0.403 & +0.504 & +0.005 \\
LC1 axial growth & $-$0.950 & +1.192 & +0.022 \\
LC2 max von Mises & $-$0.871 & +1.111 & +0.022 \\
LC2 end reaction & $-$0.886 & +1.113 & +1.354 \\
$\lambda_1$ & +0.938 & $-$1.202 & +0.772 \\
$P_{cr}$ & +0.061 & $-$0.077 & +2.146 \\
\bottomrule
\end{tabular}
\par\smallskip\parbox{\linewidth}{\footnotesize\raggedright Source: \texttt{10\_Parametric\_Study/Results/PARAMETRIC\_TRENDS.md} (central differences over the low/high cases). Families are compared only through $S$, because the step sizes differ ($\pm$10\,\% for $V$ and $q''$, $\pm$20\,\% for $t$).}
\end{table}


\section{Limitations of One-Factor-at-a-Time Analysis}
The design varies one parameter at a time about P00. It cannot identify interactions: the screening estimated the velocity $\times$ heat-flux
interaction at about 8--11\,\% of the main effects, so combined off-baseline states (for example low velocity and high heat flux together) cannot
be obtained by adding the one-factor changes. The mid-point deviation from linearity describes curvature along one axis only. No response
surface was fitted and nothing is extrapolated beyond the $\pm$10\,\% ($V$, $q''$) and $\pm$20\,\% ($t$) steps.
